\documentclass[12pt]{article}
\usepackage[T1]{fontenc}
\usepackage[utf8]{inputenc}
\usepackage{lmodern}
\usepackage{textcomp}
\usepackage{enumitem}
\usepackage{geometry}
\geometry{margin=1in}
\begin{document}
\begin{center}
Answer Key - History - K-12 Level Assessment 8 \
\small (Answers are ordered exactly as in the question paper.)
\end{center}

\section*{Multiple Choice}
\begin{enumerate}[label=\arabic*.]
\item \textbf{Answer:} A
\\ \textit{Options:} A) potlatch.; B) kiva parties.; C) totem feasts.; D) sinagua.
\\ \textit{Note:} A potlatch is a ceremonial event practiced by Indigenous peoples of the Pacific Northwest, where a host demonstrates their wealth and social status by giving away or destroying valuable items, thereby redistributing wealth within the community.
\item \textbf{Answer:} A
\\ \textit{Options:} A) the faunal assemblage from a site.; B) an osteological comparative collection from a site.; C) pumice and evidence of pyroclastic surges.; D) all of the above.
\\ \textit{Note:} The faunal assemblage from a site provides crucial information about the types of animals that were hunted or domesticated by a community, revealing their subsistence practices and dietary habits.
\item \textbf{Answer:} B
\\ \textit{Options:} A) primatology; B) archaeology; C) paleoanthropology; D) linguistics
\\ \textit{Note:} Archaeology is the subfield of anthropology that investigates past human behaviors and cultures by examining material remains, including artifacts and waste, often referred to as "garbage," to gain insights into daily life and societal practices.
\item \textbf{Answer:} D
\\ \textit{Options:} A) drinking maize beer and feasting on the hearts of slain enemies; B) sacrificing virgins in order to appease the serpent god; C) sacrificing captives of war in order to propitiate the sun god; D) sacrificing children in order to propitiate the spirits of the mountains
\\ \textit{Note:} The correct answer is based on historical accounts that indicate Capacocha was an Incan sacrificial ritual aimed at appeasing mountain deities, specifically by offering children as sacrifices.
\item \textbf{Answer:} D
\\ \textit{Options:} A) Homo erectus; B) premodern Homo sapiens; C) Neandertals; D) modern Homo sapiens
\\ \textit{Note:} Modern Homo sapiens have a proportionally larger brain size compared to their body size than other species, reflecting advanced cognitive abilities and complex social behaviors that are characteristic of humans.
\end{enumerate}

\section*{True / False}
\begin{enumerate}[label=\arabic*.]
\setcounter{enumi}{5}
\item \textbf{Answer:} True
\\ \textit{Options:} A) True; B) False
\\ \textit{Note:} The city of Uruk, along with the cultures of the Maya and the Moche, experienced significant societal challenges and declines due to periods of drought that severely impacted their agriculture and water supply, leading to food shortages and population stress.
\item \textbf{Answer:} False
\\ \textit{Options:} A) True; B) False
\\ \textit{Note:} The chasm between New Guinea/Australia and Java/Borneo is actually known as the Sunda Trench, while Beringia refers to the land bridge that once connected Asia and North America during the last Ice Age.
\item \textbf{Answer:} True
\\ \textit{Options:} A) True; B) False
\\ \textit{Note:} Both Homo sapiens and Australopithecus afarensis exhibit bipedalism, which is the ability to walk upright on two legs, indicating a shared mode of locomotion.
\item \textbf{Answer:} False
\\ \textit{Options:} A) True; B) False
\\ \textit{Note:} Uniformitarianism is a geological principle that suggests that the processes shaping the Earth today are similar to those that shaped it in the past, rather than a framework for categorizing human culture based on technology.
\item \textbf{Answer:} True
\\ \textit{Options:} A) True; B) False
\\ \textit{Note:} Homo erectus is recognized as the first hominid to exhibit advanced tool-making skills and is believed to have emerged around 1.8 million years ago, marking a significant evolutionary development in human history.
\end{enumerate}

\section*{Long Form Response}
\begin{enumerate}[label=\arabic*.]
\setcounter{enumi}{10}
\item \textbf{Answer:} The atlatl, a spear-throwing device, played a significant role in Aztec culture and its historical development, enhancing hunting techniques and efficiency. Originating in Mesoamerica, the atlatl allowed hunters to throw spears with greater force and accuracy than by hand alone, which was crucial for hunting large game such as deer and rabbits. This innovation not only improved food procurement but also contributed to the social and economic structures of the Aztec civilization, as successful hunts were vital for sustenance and trade. The atlatl thus represents a key technological advancement that shaped hunting practices and supported the growth of the Aztec society.
\\ \textit{Note:} correct because it highlights the atlatl's origins in Mesoamerica, its enhancement of hunting efficiency and accuracy, and its broader implications for the social and economic structures of Aztec civilization.
\item \textbf{Answer:} New Zealand was one of the last regions to be inhabited by humans due to its geographic isolation, located over 2,000 kilometers from the nearest landmass, which made it difficult for early humans to reach. Additionally, the harsh climate and rugged terrain posed challenges for potential settlers. The Maori, who are the indigenous people of New Zealand, arrived around the 13 th century, bringing with them rich cultural traditions and practices that evolved in isolation. This delayed occupation led to unique cultural developments, as the Maori adapted to their environment and formed a distinct identity, with strong communal ties and a deep connection to the land. The isolation also meant that when Europeans eventually arrived in the 18 th century, the Maori had developed a unique society, which had significant implications for their interactions with outsiders and their subsequent history.
\\ \textit{Note:} correct because it identifies geographic isolation and environmental challenges as key factors delaying human settlement in New Zealand, while also recognizing the significant cultural impact of the Maori's arrival and subsequent evolution in an isolated context.
\end{enumerate}

\vfill
\noindent\textit{End of Answer Key}
\end{document}